% Partie : set-theory
% Chapitre : ordinals
% Section : replacement

\documentclass[../../../include/open-logic-section]{subfiles}

\begin{document}

\olfileid[fr]{sth}{ordinals}{replacement}
\olsection{Remplacement}

Dans la \olref[sth][ordinals][vn]{sec}, nous avons motivé l'introduction
des ordinaux en proposant de les traiter comme des types d'ordre,
c'est-à-dire comme des représentants canoniques des bons ordres. Pour
que cette idée fonctionne, il faudrait démontrer que \emph{tout bon
ordre est isomorphe à un ordinal}. Nous pourrions alors définir
$\ordtype{A, <}$ comme l'ordinal $\alpha$ tel que
$\tuple{A, <} \isomorphic \alpha$.

Malheureusement, nous ne \emph{pouvons pas} démontrer ce résultat à
partir des seuls axiomes introduits jusqu'ici. (Nous verrons pourquoi
dans la \olref[replacement][strength]{sec} ; pour l'instant, retenons
simplement que nous ne le pouvons pas.) Il nous faut une nouvelle
idée, que voici :

\begin{axiom}[Schéma de remplacement]
Pour toute formule $\phi(x, y)$, l'énoncé suivant est un axiome :
\begin{quote}
	pour tout $A$, si $(\forall x \in A)\lexists![y][\phi(x,y)]$,
	alors $\Setabs{y}{(\exists x \in A)\phi(x,y)}$ existe.
\end{quote}
\end{axiom}
\noindent
Comme pour la séparation, il s'agit d'un schéma : il fournit une
infinité d'axiomes, un pour chacune des différentes formules
$\phi$, qui sont en nombre infini. On peut aussi l'écrire, comme on
le fait habituellement, sous la forme suivante :

\begin{defish}
Pour toute formule $\phi(x,y)$ qui ne contient pas «~$B$~»,
l'énoncé suivant est un axiome :
\[
\forall A[(\forall x \in A)\lexists![y][\phi(x,y)] \lif \exists B\forall y (y \in B \liff (\exists x \in A)\phi(x,y))]
\]
\end{defish}

À première vue, cette formule est toutefois assez complexe. La
conséquence immédiate du remplacement qui suit exprime peut-être
\emph{plus clairement} l'idée intuitive en jeu :\footnote{Un terme
est une expression qui désigne exactement un objet, quelle que soit
l'attribution de valeurs à ses variables libres. Par exemple,
«~$\emptyset$~» est un terme qui désigne l'ensemble vide ;
«~$\{x\}$~» est un terme qui désigne le singleton de $x$, quel que
soit $x$ ; «~$x \cup y$~» est un terme qui désigne l'union de $x$ et de
$y$, quels qu'ils soient.}

\begin{cor}
Pour tout terme $\tau(x)$ et tout ensemble $A$, l'ensemble suivant
existe :
\[
	\Setabs{\tau(x)}{x \in A} = \Setabs{y}{(\exists x \in A)y = \tau(x)}.
\]
\end{cor}

\begin{proof}
Puisque $\tau$ est un \emph{terme},
$\forall x \lexists![y][\tau(x) = y]$. A fortiori,
$(\forall x \in A)\lexists![y][\tau(x) = y]$. Par remplacement,
$\Setabs{y}{(\exists x \in A)\tau(x) = y}$ existe.
\end{proof}
\noindent
Cela suggère que «~remplacement~» est un nom approprié pour cet
axiome : étant donné un ensemble $A$, on peut former un nouvel
ensemble, $\Setabs{\tau(x)}{x \in A}$, en remplaçant chaque élément
de $A$ par son image par~$\tau$. À l'image de la notation pour
l'image d'un ensemble par une fonction, nous pourrions écrire
$\funimage{\tau}{A}$ pour $\Setabs{\tau(x)}{x \in A}$.

Il est toutefois essentiel que $\tau$ soit un \emph{terme}. Il
n'est pas nécessaire qu'il soit, ni qu'il désigne, une
\emph{fonction} au sens de la \olref[sfr][fun][rel]{sec}, c'est-à-dire
un certain ensemble de couples ordonnés. En effet, si $f$ est une
fonction en ce sens, l'ensemble
$\funimage{f}{A} = \Setabs{f(x)}{x \in A}$ n'est qu'un sous-ensemble
particulier de $\ran{f}$ ; son existence est déjà garantie par les
seuls axiomes de~$\Zminus$.\footnote{Il suffit de considérer
$\Setabs{y \in \bigcup \bigcup f}{(\exists x \in A)y = f(x)}$.}
En revanche, le remplacement est un ajout \emph{puissant} à nos
axiomes, comme nous le verrons dans le
\olref[sth][replacement][]{chap}.

\end{document}
